\documentclass[11pt,a4paper]{article}
\usepackage[T1]{fontenc}
\usepackage[utf8]{inputenc}
\usepackage{amsmath,amssymb,xcolor}
\usepackage[margin=25mm]{geometry}
\usepackage[colorlinks=true,urlcolor=blue]{hyperref}
\definecolor{nexorablue}{HTML}{2459B5}
\setlength{\parindent}{0pt}
\setlength{\parskip}{7pt}
\setlength{\emergencystretch}{2em}
\renewcommand{\familydefault}{\sfdefault}
\begin{document}

{\small\bfseries\color{nexorablue} NEXORA CPGE}\par
{\small MPSI / NEX-CIR-05 / Énoncé}\par
{\LARGE\bfseries Influence du voltmètre}\par
\vspace{0.25em}\hrule\vspace{0.6em}

Un point $B$ est relié par $60\,\mathrm{k\Omega}$ à un point de potentiel $9\,\mathrm V$, et par $30\,\mathrm{k\Omega}$ à la masse. Calculer la tension sans appareil. On mesure ensuite $V_B$ avec un voltmètre de résistance d’entrée $90\,\mathrm{k\Omega}$. Prévoir sa lecture et l’écart relatif signé par rapport à la tension sans appareil.


\end{document}
